\section{Girard's Embedding}

Here we present Girard's embedding, which is a method to represent a logical system in a different framework, often used in the context of linear logic and proof theory.\section{\texorpdfstring{Embedding of \boldmath${\la_n}$ into \boldmath${\la_b}$}{Embedding of call-by-name lambda calculus into the call-by-box lambda calculus}}

% Research more about Girard's tranlation

Let $\Lambda_{CBN}$ be the set containing all \lterms in the \lan
and let $\Lambda_{CBB}$ be the set containing all \lterms in the \lab. 
The set $\T_{CBN}$ contains all types of the \lan and $\T_{CBB}$ contains all types of the \lab. 
Now we define functions from $\Lambda_{CBN}$ to $\Lambda_{CBB}$ for term translation and from $\T_{CBN}$ to $\T_{CBB}$ for type translation. 
These functions represent an embedding of the \lan into the \lab. 

\begin{figure}[h]
  \centering
  \begin{align*}
    \girard{X}                 & = X                                          & \girard{x}             & = \boxe{x} \\
    \girard{(A_1 \fa A_2)}     & = \boxt{\girard{A_1}} \fa \girard{A_2}       & \girard{(\lamb{x}{M})} & = \lamb{x}{\girard{M}} \\
                               &                                              & \girard{(M N)}        & = \girard{M} \boxi{\girard{N}}
  \end{align*}
  \caption{Girard's translation from \lan to \lab.}
  \label{fig:girard-translation}
\end{figure}


The image of Girard's translation is a subset of the \lab given by the following grammar:

\vspace{10pt}
\begin{grammar}{
		\pr{$M, N$}{$\boxe{x} \gors \lambda x.M \gors M \boxi{N}$}
	}
\end{grammar}
\vspace{10pt}

Since application in Girard's translation is of the form 
\[
M \boxi{N}
\]

reducing the left-hand side of the application is not possible without the $\zeta$ rule that allows reduction inside boxed terms.

To clarify, we give the embedding of $\bi x$ into \lab after which we apply $\redbn$. 
It is clear that the result obtained after $\redbn$ is equal to the embedding of the result of applying $\redbn$ on $\bi x$ in the \lan, which is $\Tcbn(x) = \boxe{x}$.  

\begin{eqnarray*}
  \girard{(\bi x)} &=& \girard{(\lamb{x}{x})} \boxi{\girard{x}} \\
  &=& (\lamb{x}{\girard{x}}) \boxi{\girard{x}} \\
  &=& (\lamb{x}{\boxe{x}}) \boxi{\boxe{x}} \\
  &\redbn& (\lamb{x}{\boxe{x}}) \subst{\boxe{x}}{\boxe{x}} \\
  &=& \boxe{x}
\end{eqnarray*}

